\documentclass{ximera}

\input{../../preamble.tex}

\title{Logarithmic Analysis}

\begin{document}

\begin{abstract}
%Stuff can go here later if we want!
\end{abstract}
\maketitle





\begin{center}
\textbf{\textcolor{red!70!black}{$log_A(B)$ is the thing you raise $A$ to, to get $B$.}}
\end{center}






\begin{center}
\textbf{\textcolor{blue!55!black}{That's all you need..}}
\end{center}





\textbf{\textcolor{red!70!darkgray}{$\blacktriangleright$}} That sentence leads to all logarithmic function properties. 

\textbf{\textcolor{red!70!darkgray}{$\blacktriangleright$}} That sentence tells you how to solve all logarithmic equations. 

\textbf{\textcolor{red!70!darkgray}{$\blacktriangleright$}} That sentence reveals all of the algebra rules for logarithms. 





Our focus is function analysis.





\begin{itemize}
     \item \textbf{\textcolor{red!70!black}{Domain}} 
     \item \textbf{\textcolor{red!70!black}{Zeros}} 
     \item \textbf{\textcolor{red!70!black}{Continuity}} \\
     \text{    }$\circ$ \textbf{\textcolor{purple!85!blue}{discontinuities}}\\
     \text{    }$\circ$ \textbf{\textcolor{purple!85!blue}{singularities}}\\
     \item \textbf{\textcolor{red!70!black}{End-Behavior}} 
     \item \textbf{\textcolor{red!70!black}{Behavior}} \\
     \text{    }$\circ$ \textbf{\textcolor{purple!85!blue}{intervals where increasing}}\\
     \text{    }$\circ$ \textbf{\textcolor{purple!85!blue}{intervals where decreasing}}\\
     \item \textbf{\textcolor{red!70!black}{Global Maximum and Minimum}} 
     \item \textbf{\textcolor{red!70!black}{Local Maximums and Minimums}} 
     \item \textbf{\textcolor{red!70!black}{Range}} 
     \item \textbf{\textcolor{blue!55!black}{...and we would like a nice graph}} 
\end{itemize}




\subsection*{Learning Outcomes}


\begin{sectionOutcomes}
In this section, students will 

\begin{itemize}
\item paint a full picture of logarithmic functions.
\end{itemize}
\end{sectionOutcomes}













\begin{center}
\textbf{\textcolor{green!50!black}{ooooo-=-=-=-ooOoo-=-=-=-ooooo}} \\

more examples can be found by following this link\\ \link[More Examples of Logarithmic Functions]{https://ximera.osu.edu/csccmathematics/precalculus2/precalculus2/logFunctions/examples/exampleList}

\end{center}










\end{document}
